\documentclass{ximera}

\input{../../../preamble.tex}


\definecolor{fvdnavy}{HTML}{071D43}
\definecolor{fvdcyan}{HTML}{20BFC6}
\definecolor{fvdcyanlight}{HTML}{EFFCFB}
\definecolor{fvdmagenta}{HTML}{F10D69}
\definecolor{fvdmagentalight}{HTML}{FFF1F7}
\definecolor{fvdtext}{HTML}{163044}





%=================================================
% Pedagogische omgevingen
% PDF: tcolorbox
% HTML: learning-box
%=================================================

\newenvironment{voorkennis}
{%
    \ifdefined\HCode
        \HCode{%
<section class="learning-box learning-box--prior">
<div class="learning-box__header">
<span class="learning-box__icon" aria-hidden="true"></span>
<span class="learning-box__title">Voorkennis</span>
</div>
<div class="learning-box__content">
        }%
        \begin{itemize}
    \else
       \begin{tcolorbox}[
    enhanced,
    breakable,
    colback=fvdcyanlight,
    colframe=fvdcyan,
    coltext=fvdtext,
    boxrule=0.5pt,
    leftrule=4pt,
    arc=4mm,
    outer arc=4mm,
    left=7mm,
    right=6mm,
    top=5mm,
    bottom=5mm,
    before skip=6mm,
    after skip=6mm,
    title=Voorkennis,
    coltitle=fvdcyan,
    fonttitle=\bfseries\Large,
    attach title to upper={\par\smallskip},
    boxed title style={
        empty,
        left=0mm,
        right=0mm,
        top=0mm,
        bottom=0mm
    },
    overlay={
        \draw[
            fvdcyan,
            line width=0.5pt
        ]
        ([xshift=42mm,yshift=-13mm]frame.north west)
        --
        ([xshift=-7mm,yshift=-13mm]frame.north east);
    }
]
        \begin{itemize}
    \fi
}
{%
    \end{itemize}
    \ifdefined\HCode
        \HCode{%
</div>
</section>
        }%
    \else
        \end{tcolorbox}
    \fi
}


\newenvironment{leerdoelen}
{%
    \ifdefined\HCode
        \HCode{%
<section class="learning-box learning-box--goals">
<div class="learning-box__header">
<span class="learning-box__icon" aria-hidden="true"></span>
<span class="learning-box__title">Leerdoelen</span>
</div>
<div class="learning-box__content">
        }%
        \begin{itemize}
    \else
\begin{tcolorbox}[
    enhanced,
    breakable,
    colback=fvdmagentalight,
    colframe=fvdmagenta,
    coltext=fvdtext,
    boxrule=0.5pt,
    leftrule=4pt,
    arc=4mm,
    outer arc=4mm,
    left=7mm,
    right=6mm,
    top=5mm,
    bottom=5mm,
    before skip=6mm,
    after skip=6mm,
    title=Leerdoelen,
    coltitle=fvdmagenta,
    fonttitle=\bfseries\Large,
    attach title to upper={\par\smallskip},
    boxed title style={
        empty,
        left=0mm,
        right=0mm,
        top=0mm,
        bottom=0mm
    },
    overlay={
        \draw[
            fvdmagenta,
            line width=0.5pt
        ]
        ([xshift=42mm,yshift=-13mm]frame.north west)
        --
        ([xshift=-7mm,yshift=-13mm]frame.north east);
    }
]
        \begin{itemize}
    \fi
}
{%
    \end{itemize}
    \ifdefined\HCode
        \HCode{%
</div>
</section>
        }%
    \else
        \end{tcolorbox}
    \fi
}


\addPrintStyle{../../..}

\title{De gereedschapskist van de fysicus}
\author{Ingmar Herreman}

\begin{abstract}
Fysica begint met waarnemen en meten. Een afstand, een tijdsduur,
een massa of een temperatuur wordt pas bruikbare wetenschappelijke
informatie wanneer duidelijk is welke grootheid gemeten werd, in welke
eenheid het resultaat wordt uitgedrukt en hoe nauwkeurig de meting is.

In dit eerste hoofdstuk bouwen we daarom een gereedschapskist op die
we doorheen de volledige fysicacursus zullen gebruiken. We werken met
fysische grootheden en SI-eenheden, wetenschappelijke notatie,
voorvoegsels, meetnauwkeurigheid en beduidende cijfers.

Daarna maken we de stap van scalaire naar vectoriële grootheden.
Krachten, snelheden en versnellingen hebben immers niet alleen een
grootte, maar ook een richting en een zin. We leren vectoren voorstellen,
optellen en ontbinden in componenten.

Ook experimenteren krijgt een plaats. We onderzoeken hoe een
wetenschappelijke onderzoeksvraag wordt vertaald naar meetbare
grootheden en hoe meetresultaten kritisch worden verwerkt en
geïnterpreteerd.

Ten slotte ontwikkelen we een vaste aanpak voor fysische problemen:
de situatie begrijpen, een geschikt model kiezen, systematisch rekenen
en het resultaat kritisch controleren.

Deze gereedschappen vormen de basis voor het volgende hoofdstuk,
waarin we krachten en de wetten van Newton bestuderen.
\end{abstract}

\begin{document}

% FVD_CHAPTER_DATA_START
% Automatisch gegenereerd uit index.tex.
% Niet handmatig aanpassen.
\fvdchapterdata{1}{Gereedschapskist en Newton}{1}
% FVD_CHAPTER_DATA_END

\maketitle


% ============================================================
% VOORKENNIS EN LEERDOELEN
% ============================================================

\begin{voorkennis}
  \item Je kan eenvoudige meetresultaten met een getalwaarde en eenheid noteren.
  \item Je kan rekenen met machten van tien.
  \item Je kent enkele courante SI-eenheden en voorvoegsels.
  \item Je kan eenvoudige formules gebruiken en een onbekende grootheid vrijmaken.
  \item Je kent uit de tweede graad het onderscheid tussen scalaire en vectoriële grootheden.
  \item Je kan de stelling van Pythagoras en eenvoudige goniometrie gebruiken.
\end{voorkennis}

\begin{leerdoelen}
  \item Je kan fysische grootheden, grootheidssymbolen, getalwaarden en eenheden correct onderscheiden.
  \item Je kan met SI-eenheden en courante voorvoegsels werken.
  \item Je kan eenheden omzetten en daarbij ook machten van eenheden correct behandelen.
  \item Je kan getallen schrijven en interpreteren in wetenschappelijke notatie.
  \item Je kan meetresultaten met een verantwoord aantal beduidende cijfers rapporteren.
  \item Je kan het onderscheid uitleggen tussen resolutie, meetnauwkeurigheid en beduidende cijfers.
  \item Je kan scalaire en vectoriële grootheden onderscheiden.
  \item Je kan de grootte, richting en zin van een vector beschrijven.
  \item Je kan vectoren grafisch optellen en ontbinden in componenten.
  \item Je kan uit vectorcomponenten de grootte en richting van een vector bepalen.
  \item Je kan bij een eenvoudig experiment onafhankelijke, afhankelijke en controlevariabelen onderscheiden.
  \item Je kan meetgegevens kritisch beoordelen.
  \item Je kan een fysische berekening controleren met behulp van eenheden.
  \item Je kan een fysisch probleem systematisch mathematiseren, oplossen en interpreteren.
\end{leerdoelen}


% ============================================================
% 1 WAAROM?
% ============================================================

\FVDsection{Waarom heb je een gereedschapskist nodig?}

Een afstand van $12{,}4\,\mathrm m$, een tijd van $3{,}8\,\mathrm s$,
een snelheid van $90\,\mathrm{km/h}$ of een kracht van
$250\,\mathrm N$: fysica zit vol getallen.

Maar fysica is niet simpelweg rekenen met getallen.

Bij elk getal moeten we ons afvragen:

\begin{itemize}
  \item Welke fysische grootheid wordt beschreven?
  \item Welke eenheid wordt gebruikt?
  \item Hoe nauwkeurig is het resultaat gekend?
  \item Is de gekozen notatie geschikt?
  \item Heeft de grootheid ook een richting?
  \item Past het resultaat bij de werkelijkheid?
\end{itemize}

De fysicus probeert waarnemingen om te zetten in een beschrijving die
kan worden onderzocht, gecontroleerd en gebruikt om voorspellingen te
maken.

Een mogelijke keten is:

\[
\boxed{
\text{waarnemen}
\longrightarrow
\text{meten}
\longrightarrow
\text{modelleren}
\longrightarrow
\text{verklaren}
\longrightarrow
\text{voorspellen}
}
\]

De begrippen en technieken uit dit hoofdstuk zullen daarom niet alleen
hier gebruikt worden. Ze keren in bijna elk volgend hoofdstuk terug.


% ============================================================
% 2 FYSISCHE GROOTHEDEN
% ============================================================

\FVDsection{Fysische grootheden en eenheden}

Een eigenschap die we kunnen meten of berekenen noemen we een
\textbf{fysische grootheid}.

Voorbeelden zijn:

\begin{itemize}
  \item lengte;
  \item tijd;
  \item massa;
  \item temperatuur;
  \item snelheid;
  \item versnelling;
  \item kracht;
  \item energie.
\end{itemize}

\begin{definitie}[Fysische grootheid]

Een \textbf{fysische grootheid} is een meetbare of berekenbare
eigenschap van een voorwerp, systeem of verschijnsel.

\end{definitie}

Iedere grootheid krijgt meestal een symbool.

Bijvoorbeeld:

\[
\begin{array}{c|c}
\text{grootheid} & \text{symbool}\\
\hline
\text{tijd} & t\\
\text{massa} & m\\
\text{lengte} & l\\
\text{snelheid} & v\\
\text{versnelling} & a\\
\text{kracht} & F\\
\text{energie} & E
\end{array}
\]

Een gemeten waarde bestaat uit een getalwaarde en een eenheid.

Bijvoorbeeld:

\[
l=2{,}35\,\mathrm m.
\]

Daarin is:

\[
\boxed{
\begin{array}{ll}
l & \text{het grootheidssymbool},\\
2{,}35 & \text{de getalwaarde},\\
\mathrm m & \text{de eenheid}.
\end{array}
}
\]

\begin{opmerking}

De uitspraak

\[
d=12
\]

is in een fysische context meestal onvolledig.

De uitspraak

\[
d=12\,\mathrm m
\]

bevat wel de noodzakelijke informatie over de gebruikte eenheid.

\end{opmerking}


% ============================================================
% 3 SI
% ============================================================

\FVDsection{Het SI-stelsel}

Wetenschappers moeten meetresultaten met elkaar kunnen vergelijken.

Daarom wordt internationaal gewerkt met het
\textbf{Internationaal Stelsel van Eenheden}, kortweg het
\textbf{SI-stelsel}.

Enkele belangrijke SI-basiseenheden zijn:

\[
\begin{array}{c|c|c}
\text{grootheid} & \text{symbool} & \text{SI-eenheid}\\
\hline
\text{lengte} & l & \mathrm m\\
\text{massa} & m & \mathrm{kg}\\
\text{tijd} & t & \mathrm s\\
\text{temperatuur} & T & \mathrm K\\
\text{stroomsterkte} & I & \mathrm A\\
\text{hoeveelheid stof} & n & \mathrm{mol}\\
\text{lichtsterkte} & I_v & \mathrm{cd}
\end{array}
\]

Veel andere fysische eenheden zijn afgeleid van deze basiseenheden.

Voor snelheid:

\[
v=\frac{\Delta x}{\Delta t}.
\]

De eenheid wordt:

\[
[v]
=
\frac{\mathrm m}{\mathrm s}
=
\mathrm{m/s}.
\]

Voor versnelling:

\[
[a]
=
\frac{\mathrm{m/s}}{\mathrm s}
=
\mathrm{m/s^2}.
\]

Later zullen we zien dat voor kracht geldt:

\[
1\,\mathrm N
=
1\,\mathrm{kg\,m/s^2}.
\]


\begin{oefening}[Grootheid of eenheid?]{\oefB\ESS}

Geef telkens de grootheid, het grootheidssymbool en de eenheid.

\begin{enumerate}
  \item Een experiment duurt $12{,}5\,\mathrm s$.
  \item Een voorwerp heeft een massa van $2{,}40\,\mathrm{kg}$.
  \item Een wagen rijdt $25\,\mathrm{m/s}$.
  \item Op een voorwerp werkt een kracht van $18\,\mathrm N$.
  \item Een toestel verbruikt $850\,\mathrm J$ energie.
\end{enumerate}

\end{oefening}


% ============================================================
% 4 VOORVOEGSELS
% ============================================================

\FVDsection{Van nanometer tot gigameter}

In de fysica komen zeer verschillende lengteschalen voor.

Een atoom heeft een typische grootteorde van ongeveer

\[
10^{-10}\,\mathrm m,
\]

terwijl afstanden in de astronomie vele miljarden meters kunnen
bedragen.

Daarom gebruiken we \textbf{voorvoegsels}.

\[
\begin{array}{c|c|c}
\text{voorvoegsel} & \text{symbool} & \text{factor}\\
\hline
\text{tera}  & \mathrm T & 10^{12}\\
\text{giga}  & \mathrm G & 10^9\\
\text{mega}  & \mathrm M & 10^6\\
\text{kilo}  & \mathrm k & 10^3\\
\text{deci}  & \mathrm d & 10^{-1}\\
\text{centi} & \mathrm c & 10^{-2}\\
\text{milli} & \mathrm m & 10^{-3}\\
\text{micro} & \mu & 10^{-6}\\
\text{nano}  & \mathrm n & 10^{-9}\\
\text{pico}  & \mathrm p & 10^{-12}
\end{array}
\]


\FVDsubsection{Eenheden omzetten}

Beschouw:

\[
3{,}5\,\mathrm{km}.
\]

Omdat:

\[
1\,\mathrm{km}=10^3\,\mathrm m,
\]

geldt:

\[
3{,}5\,\mathrm{km}
=
3{,}5\cdot10^3\,\mathrm m
=
3500\,\mathrm m.
\]


\begin{voorbeeld}[Van kilometer per uur naar meter per seconde]

Een auto rijdt:

\[
90\,\mathrm{km/h}.
\]

We gebruiken:

\[
1\,\mathrm{km}=1000\,\mathrm m
\]

en:

\[
1\,\mathrm h=3600\,\mathrm s.
\]

Dan:

\[
90\,\mathrm{km/h}
=
90\cdot
\frac{1000\,\mathrm m}{3600\,\mathrm s}.
\]

Dus:

\[
\boxed{
v=25\,\mathrm{m/s}.
}
\]

\end{voorbeeld}


\FVDsubsection{Opgelet met oppervlakte en volume}

Bij oppervlakte-eenheden moet ook de omzettingsfactor gekwadrateerd
worden.

Omdat:

\[
1\,\mathrm{cm}
=
10^{-2}\,\mathrm m,
\]

geldt:

\[
1\,\mathrm{cm^2}
=
(10^{-2}\,\mathrm m)^2.
\]

Dus:

\[
\boxed{
1\,\mathrm{cm^2}
=
10^{-4}\,\mathrm{m^2}.
}
\]

Voor volume:

\[
1\,\mathrm{cm^3}
=
(10^{-2}\,\mathrm m)^3
=
10^{-6}\,\mathrm{m^3}.
\]


\begin{opmerking}

De omzetting

\[
1\,\mathrm{cm^2}
=
10^{-2}\,\mathrm{m^2}
\]

is fout.

Het voorvoegsel hoort bij de volledige lengteeenheid en wordt dus mee
gekwadrateerd.

\end{opmerking}


\begin{oefening}[Eenheden omzetten]{\oefB\ESS}

Zet om.

\begin{enumerate}
  \item $4{,}2\,\mathrm{km}$ naar meter.
  \item $350\,\mathrm{ms}$ naar seconde.
  \item $7{,}4\,\mathrm{cm}$ naar meter.
  \item $72\,\mathrm{km/h}$ naar $\mathrm{m/s}$.
  \item $5{,}0\,\mathrm{cm^2}$ naar $\mathrm{m^2}$.
  \item $12\,\mathrm{cm^3}$ naar $\mathrm{m^3}$.
\end{enumerate}

\end{oefening}


% ============================================================
% 5 WETENSCHAPPELIJKE NOTATIE
% ============================================================

\FVDsection{Wetenschappelijke notatie}

Wetenschappelijke notatie maakt zeer grote en zeer kleine getallen
overzichtelijk.

De lichtsnelheid is ongeveer:

\[
300\,000\,000\,\mathrm{m/s}.
\]

Dat schrijven we als:

\[
3{,}00\cdot10^8\,\mathrm{m/s}.
\]

De grootteorde van een atoom is ongeveer:

\[
0{,}0000000001\,\mathrm m.
\]

Dat wordt:

\[
1\cdot10^{-10}\,\mathrm m.
\]


\begin{definitie}[Wetenschappelijke notatie]

Een getal staat in \textbf{wetenschappelijke notatie} wanneer het
geschreven wordt als:

\[
\boxed{
a\cdot10^n
}
\]

met:

\[
1\leq |a|<10
\]

en:

\[
n\in\mathbb Z.
\]

\end{definitie}


\begin{voorbeeld}

Schrijf:

\[
0{,}000061
\]

in wetenschappelijke notatie.

We verschuiven de komma vijf plaatsen naar rechts:

\[
0{,}000061
=
6{,}1\cdot10^{-5}.
\]

Dus:

\[
\boxed{
6{,}1\cdot10^{-5}
}
\]

\end{voorbeeld}


\begin{oefening}[Wetenschappelijke notatie]{\oefB\ESS}

Schrijf in wetenschappelijke notatie.

\begin{enumerate}
  \item $470000$
  \item $0{,}000832$
  \item $-0{,}00000725$
  \item $8600000000$
  \item $0{,}00000000045$
\end{enumerate}

Schrijf daarna opnieuw in gewone decimale notatie:

\begin{enumerate}
  \item $3{,}8\cdot10^6$
  \item $6{,}02\cdot10^{-4}$
  \item $9{,}1\cdot10^{-9}$
\end{enumerate}

\end{oefening}


% ============================================================
% 6 METEN
% ============================================================

\FVDsection{Een meting is nooit exact}

Stel dat we de lengte van een tafel meten en vinden:

\[
l=1{,}84\,\mathrm m.
\]

Betekent dit dat de tafel exact:

\[
1{,}840000000000\ldots\,\mathrm m
\]

lang is?

Nee.

Iedere meting heeft een beperkte nauwkeurigheid.

De kwaliteit van een meetresultaat hangt onder andere af van:

\begin{itemize}
  \item het gebruikte meetinstrument;
  \item de resolutie van het meetinstrument;
  \item de meetmethode;
  \item de omstandigheden;
  \item de persoon die meet;
  \item eventuele systematische afwijkingen.
\end{itemize}


\FVDsubsection{Resolutie}

\begin{definitie}[Resolutie]

De \textbf{resolutie} van een meetinstrument is de kleinste verandering
van de gemeten grootheid die het instrument kan onderscheiden.

\end{definitie}

Een gewone liniaal met millimeterverdeling heeft bijvoorbeeld een
andere resolutie dan een schuifmaat.

Een digitaal toestel kan bijvoorbeeld tonen:

\[
12{,}37\,\mathrm V.
\]

De laatste weergegeven positie geeft informatie over de resolutie van
het toestel.


\begin{opmerking}

Meer cijfers op een digitaal scherm betekenen niet automatisch dat het
meetresultaat ook werkelijk nauwkeuriger is.

Ook kalibratie, meetmethode en experimentele omstandigheden kunnen de
betrouwbaarheid beperken.

\end{opmerking}


% ============================================================
% 7 BEDUIDENDE CIJFERS
% ============================================================

\FVDsection{Beduidende cijfers}

Wanneer we een meetresultaat noteren, moeten de cijfers overeenkomen
met de nauwkeurigheid die we werkelijk kunnen verantwoorden.

Daarom gebruiken we het begrip \textbf{beduidende cijfers}.


\begin{definitie}[Beduidende cijfers]

De \textbf{beduidende cijfers} van een meetwaarde zijn de cijfers die
informatie dragen over de gerapporteerde waarde en haar nauwkeurigheid.

\end{definitie}

Beschouw:

\[
12{,}4\,\mathrm{cm}.
\]

Dit getal bevat drie beduidende cijfers:

\[
1,\qquad 2,\qquad 4.
\]


\FVDsubsection{Praktische regels}

\begin{enumerate}

  \item Alle cijfers van $1$ tot en met $9$ zijn beduidend.

  Bijvoorbeeld:

  \[
  37{,}2
  \]

  heeft drie beduidende cijfers.

  \item Nullen tussen andere beduidende cijfers zijn beduidend.

  Bijvoorbeeld:

  \[
  4{,}007
  \]

  heeft vier beduidende cijfers.

  \item Nullen vóór het eerste niet-nulcijfer zijn niet beduidend.

  Bijvoorbeeld:

  \[
  0{,}00450
  \]

  heeft drie beduidende cijfers:

  \[
  4,\quad5,\quad0.
  \]

  \item Nullen achter de komma en na een niet-nulcijfer kunnen wel
  beduidend zijn.

  Bijvoorbeeld:

  \[
  2{,}300
  \]

  heeft vier beduidende cijfers.

\end{enumerate}


\FVDsubsection{Waarom wetenschappelijke notatie helpt}

Bij een getal zoals:

\[
5000
\]

is zonder verdere context niet altijd duidelijk hoeveel beduidende
cijfers bedoeld zijn.

Wetenschappelijke notatie maakt dat ondubbelzinnig.

\[
5\cdot10^3
\]

heeft één beduidend cijfer.

\[
5{,}0\cdot10^3
\]

heeft twee beduidende cijfers.

\[
5{,}000\cdot10^3
\]

heeft vier beduidende cijfers.


\begin{opmerking}

Wetenschappelijke notatie geeft dus niet alleen de grootteorde van een
getal weer. Ze kan ook duidelijk maken met hoeveel beduidende cijfers
een meetwaarde wordt gerapporteerd.

\end{opmerking}


\begin{oefening}[Beduidende cijfers herkennen]{\oefB\ESS}

Bepaal het aantal beduidende cijfers.

\begin{enumerate}
  \item $0{,}00450$
  \item $2{,}030$
  \item $6{,}00\cdot10^4$
  \item $0{,}0000700$
  \item $1{,}0050$
\end{enumerate}

\end{oefening}


% ============================================================
% 8 REKENEN MET MEETWAARDEN
% ============================================================

\FVDsection{Rekenen met meetwaarden}

Een rekenmachine kan veel cijfers produceren.

Dat betekent niet dat al die cijfers fysisch betekenisvol zijn.


\begin{voorbeeld}[Een gemiddelde snelheid]

Een voorwerp legt:

\[
s=12{,}4\,\mathrm m
\]

af in:

\[
t=3{,}2\,\mathrm s.
\]

De rekenmachine geeft:

\[
v
=
\frac{12{,}4}{3{,}2}
=
3{,}875\,\mathrm{m/s}.
\]

De tijd is gegeven met twee beduidende cijfers.

Een verantwoord gerapporteerd resultaat is daarom:

\[
\boxed{
v=3{,}9\,\mathrm{m/s}.
}
\]

Het resultaat:

\[
3{,}875000000\,\mathrm{m/s}
\]

suggereert een nauwkeurigheid die de oorspronkelijke metingen niet
hebben.

\end{voorbeeld}


\FVDsubsection{Vermenigvuldigen en delen}

Bij vermenigvuldigen en delen wordt een resultaat in deze cursus
doorgaans gerapporteerd met evenveel beduidende cijfers als de
meetwaarde met het kleinste aantal beduidende cijfers.


\begin{voorbeeld}

Bereken:

\[
(2{,}30\cdot10^2)
(1{,}0\cdot10^{-3}).
\]

We vinden:

\[
0{,}230.
\]

De factor:

\[
1{,}0\cdot10^{-3}
\]

heeft twee beduidende cijfers.

Daarom rapporteren we:

\[
\boxed{
0{,}23.
}
\]

\end{voorbeeld}


\FVDsubsection{Optellen en aftrekken}

Bij optellen en aftrekken kijken we niet naar het aantal beduidende
cijfers, maar naar de plaats van het laatste betrouwbare cijfer.

Bijvoorbeeld:

\[
12{,}35+1{,}2=13{,}55.
\]

De tweede meetwaarde is slechts tot op één decimaal gegeven.

Daarom:

\[
\boxed{
13{,}6.
}
\]


\begin{oefening}[Te veel cijfers]{\oefU\ESS}

Een leerling meet:

\[
l=12{,}4\,\mathrm{cm}
\]

en:

\[
b=3{,}2\,\mathrm{cm}.
\]

Hij berekent:

\[
A=39{,}680000\,\mathrm{cm^2}.
\]

\begin{question}
Leg uit waarom deze notatie fysisch misleidend is.
\end{question}

\begin{question}
Bereken een verantwoord gerapporteerde waarde van de oppervlakte.
\end{question}

\begin{question}
Schrijf het resultaat ook in wetenschappelijke notatie zodat het aantal
beduidende cijfers ondubbelzinnig is.
\end{question}

\end{oefening}


\begin{oefening}[Een meetresultaat beoordelen]{\oefU\ESS}

Een leerling schrijft in een practicumverslag:

\begin{quote}
De lengte van een metalen staaf is
$12{,}3478291\,\mathrm{cm}$.
De lengte werd gemeten met een gewone schoolliniaal met
millimeterverdeling.
\end{quote}

Analyseer deze rapportering.

\begin{question}
Wat is problematisch aan het opgegeven aantal cijfers?
\end{question}

\begin{question}
Waarom kan een rekenmachine of digitaal document geen extra
meetnauwkeurigheid creëren?
\end{question}

\begin{question}
Hoe zou je het resultaat realistischer rapporteren?
\end{question}

\end{oefening}


% ============================================================
% 9 EXPERIMENTEREN
% ============================================================

\FVDsection{Van meten naar onderzoeken}

Meten is geen doel op zich.

Een fysicus meet om een vraag over de werkelijkheid te beantwoorden.

Bijvoorbeeld:

\begin{center}
\textit{Hoe verandert de uitrekking van een veer wanneer de kracht
op de veer toeneemt?}
\end{center}

Dat is een vraag die we experimenteel kunnen onderzoeken.


\FVDsubsection{Onderzoeksvraag}

Een goede experimentele onderzoeksvraag bevat meestal meetbare
grootheden.

Bijvoorbeeld:

\[
\boxed{
\text{Hoe hangt de uitrekking }x
\text{ af van de aangelegde kracht }F?
}
\]


\FVDsubsection{Hypothese}

\begin{definitie}[Hypothese]

Een \textbf{hypothese} is een toetsbare verwachting over het resultaat
van een onderzoek.

\end{definitie}

Bijvoorbeeld:

\begin{quote}
Als de kracht op de veer toeneemt, dan verwachten we dat de uitrekking
toeneemt.
\end{quote}


\FVDsubsection{Variabelen}

In een gecontroleerd experiment onderscheiden we verschillende soorten
variabelen.

\begin{definitie}[Onafhankelijke variabele]

De \textbf{onafhankelijke variabele} is de grootheid die de onderzoeker
doelbewust verandert.

\end{definitie}

\begin{definitie}[Afhankelijke variabele]

De \textbf{afhankelijke variabele} is de grootheid waarvan onderzocht
wordt hoe ze reageert op de verandering van de onafhankelijke
variabele.

\end{definitie}

\begin{definitie}[Controlevariabelen]

\textbf{Controlevariabelen} zijn andere relevante factoren die tijdens
het experiment zoveel mogelijk constant gehouden worden.

\end{definitie}


\begin{voorbeeld}[Een veer onderzoeken]

We onderzoeken hoe de uitrekking van één bepaalde veer afhangt van
de aangelegde kracht.

Dan kan gelden:

\[
\begin{array}{ll}
\text{onafhankelijke variabele} & F,\\
\text{afhankelijke variabele} & x,\\
\text{controlevariabelen} &
\text{dezelfde veer, meetmethode, opstelling, \ldots}
\end{array}
\]

\end{voorbeeld}


\FVDsubsection{Een mogelijke onderzoekscyclus}

Een wetenschappelijk onderzoek verloopt niet altijd exact lineair,
maar volgende stappen vormen een bruikbaar kader:

\[
\boxed{
\begin{array}{c}
\text{waarneming of probleem}\\
\downarrow\\
\text{onderzoeksvraag}\\
\downarrow\\
\text{hypothese}\\
\downarrow\\
\text{onderzoeksplan}\\
\downarrow\\
\text{metingen}\\
\downarrow\\
\text{gegevens verwerken}\\
\downarrow\\
\text{besluit}\\
\downarrow\\
\text{reflectie}
\end{array}
}
\]

Meetgegevens kunnen worden weergegeven met:

\begin{itemize}
  \item tabellen;
  \item grafieken;
  \item functievoorschriften;
  \item trendlijnen;
  \item numerieke of grafische modellen.
\end{itemize}


\FVDsubsection{Kritisch omgaan met meetgegevens}

Na een experiment vragen we niet alleen:

\begin{center}
\textit{Wat is het antwoord?}
\end{center}

We vragen ook:

\begin{itemize}
  \item Zijn er voldoende metingen uitgevoerd?
  \item Zijn de meetwaarden reproduceerbaar?
  \item Zijn er uitschieters?
  \item Is het meetinstrument voldoende nauwkeurig?
  \item Zijn relevante controlevariabelen constant gehouden?
  \item Past het gekozen model bij de meetgegevens?
  \item Welke beperkingen heeft het experiment?
\end{itemize}


\begin{opmerking}

Een afwijkend meetpunt mag niet zomaar verwijderd worden omdat het
niet mooi bij de andere meetpunten past.

Er moet eerst onderzocht worden of er een inhoudelijke reden is om
dat meetpunt als onbetrouwbaar te beschouwen.

\end{opmerking}


\begin{oefening}[Een opvallend meetpunt]{\oefU\ESS}

Een leerling meet vijf keer de tijd die een karretje nodig heeft om
dezelfde afstand af te leggen.

\[
\begin{array}{c|ccccc}
\text{meting} & 1 & 2 & 3 & 4 & 5\\
\hline
t\,(\mathrm s) & 2{,}41 & 2{,}39 & 2{,}40 & 3{,}12 & 2{,}42
\end{array}
\]

\begin{question}
Welke meetwaarde valt op?
\end{question}

\begin{question}
Mag die meetwaarde onmiddellijk worden verwijderd?
\end{question}

\begin{question}
Geef minstens twee mogelijke verklaringen voor de afwijking.
\end{question}

\begin{question}
Welke bijkomende informatie zou je willen hebben voordat je beslist
hoe je met de meetwaarde omgaat?
\end{question}

\end{oefening}


% ============================================================
% 10 SCALAIR EN VECTOR
% ============================================================

\FVDsection{Niet alles past in één getal}

Stel dat iemand zegt:

\begin{center}
``Een drone moet zich $20\,\mathrm m$ verplaatsen.''
\end{center}

Is dat voldoende informatie om de drone naar zijn bestemming te sturen?

Nee.

We weten hoe groot de verplaatsing moet zijn, maar niet in welke
richting.

De uitspraak:

\begin{center}
``$20\,\mathrm m$ naar het oosten''
\end{center}

bevat meer informatie.

Daarmee komen we bij het onderscheid tussen scalaire en vectoriële
grootheden.


\begin{definitie}[Scalaire grootheid]

Een \textbf{scalaire grootheid} wordt volledig beschreven door een
grootte, uitgedrukt met een getalwaarde en eenheid.

\end{definitie}

Voorbeelden zijn:

\[
m=5{,}0\,\mathrm{kg},
\]

\[
T=293\,\mathrm K,
\]

\[
E=120\,\mathrm J.
\]


\begin{definitie}[Vectoriële grootheid]

Een \textbf{vectoriële grootheid} heeft naast een grootte ook een
richting en een zin.

\end{definitie}

We stellen vectoriële grootheden bijvoorbeeld voor als:

\[
\vec v,
\qquad
\vec a,
\qquad
\vec F.
\]

De grootte van een vector schrijven we als:

\[
|\vec v|
\]

of, wanneer er geen verwarring mogelijk is, eenvoudig als:

\[
v.
\]


\begin{opmerking}

Bij sommige fysische vectoren, in het bijzonder krachten, kan ook het
aangrijpingspunt fysisch belangrijk zijn.

Dat aangrijpingspunt is echter geen algemene extra eigenschap van een
abstracte wiskundige vector.

\end{opmerking}


\begin{oefening}[Scalair of vectorieel?]{\oefB\ESS}

Bepaal telkens of de grootheid scalair of vectorieel is.

\begin{enumerate}
  \item massa;
  \item kracht;
  \item energie;
  \item snelheid;
  \item temperatuur;
  \item versnelling;
  \item volume;
  \item elektrisch veld;
  \item tijd;
  \item verplaatsing.
\end{enumerate}

\end{oefening}


% ============================================================
% 11 VECTOR VOORSTELLEN
% ============================================================

\FVDsection{Een vector voorstellen}

Een vector wordt grafisch voorgesteld door een pijl.

Daarbij stellen we voor:

\[
\boxed{
\text{grootte, richting en zin}
}
\]

De lengte van de pijl stelt de grootte voor.

De ligging van de pijl bepaalt de richting.

De pijlpunt bepaalt de zin.


\begin{voorbeeld}[Een snelheidsvector]

Een wagen rijdt met een snelheid van:

\[
72\,\mathrm{km/h}
\]

naar het noordoosten.

We schrijven de snelheid als:

\[
\vec v.
\]

De grootte is:

\[
|\vec v|
=
72\,\mathrm{km/h}.
\]

Om de snelheid volledig te beschrijven moet ook de richting en zin
worden aangegeven.

\end{voorbeeld}


\FVDsubsection{Schaal bij een vectortekening}

Bij een vectortekening kan een schaal worden gebruikt.

Bijvoorbeeld:

\[
1\,\mathrm{cm}
\longleftrightarrow
10\,\mathrm N.
\]

Een kracht van:

\[
40\,\mathrm N
\]

wordt dan voorgesteld door een pijl van:

\[
4\,\mathrm{cm}.
\]


% ============================================================
% 12 VECTOR OPTELLEN
% ============================================================

\FVDsection{Vectoren optellen}

Bij scalaire grootheden volstaat gewone algebra vaak.

Bij vectoren moeten ook richting en zin worden meegenomen.


\FVDsubsection{Dezelfde richting en dezelfde zin}

Twee vectoren in dezelfde richting en zin versterken elkaar.

Als:

\[
F_1=3\,\mathrm N
\]

en:

\[
F_2=5\,\mathrm N
\]

beide naar rechts wijzen, dan:

\[
F_{\mathrm{res}}
=
3+5
=
8\,\mathrm N.
\]

Dus:

\[
\boxed{
F_{\mathrm{res}}
=
8\,\mathrm N
\text{ naar rechts}.
}
\]


\FVDsubsection{Dezelfde richting, tegengestelde zin}

Stel:

\[
F_1=8\,\mathrm N
\]

naar rechts en:

\[
F_2=3\,\mathrm N
\]

naar links.

De resulterende kracht is:

\[
\boxed{
F_{\mathrm{res}}
=
5\,\mathrm N
\text{ naar rechts}.
}
\]


\FVDsubsection{De kop-staartmethode}

Voor vectoren met verschillende richtingen kunnen we de
\textbf{kop-staartmethode} gebruiken.

We verschuiven een vector evenwijdig zodat zijn beginpunt samenvalt
met het eindpunt van de vorige vector.

De resultante loopt van het beginpunt van de eerste vector naar het
eindpunt van de laatste vector.

Symbolisch:

\[
\boxed{
\vec A+\vec B=\vec R.
}
\]


\begin{oefening}[Kan je vectoren gewoon optellen?]{\oefU\ESS}

Een leerling beweert:

\begin{center}
``De grootte van $\vec A+\vec B$ is altijd gelijk aan $A+B$.''
\end{center}

Onderzoek deze uitspraak.

\begin{question}
Geef een situatie waarin de uitspraak klopt.
\end{question}

\begin{question}
Geef een situatie waarin de uitspraak niet klopt.
\end{question}

\begin{question}
Leg met een vectorschets uit waarom de richting van de vectoren
belangrijk is.
\end{question}

\end{oefening}


% ============================================================
% 13 COMPONENTEN
% ============================================================

\FVDsection{Een vector ontbinden}

Een vector kunnen we ontbinden in loodrechte componenten.

In twee dimensies:

\[
\boxed{
\vec A
=
\vec A_x+\vec A_y.
}
\]

Als een vector $\vec A$ een hoek $\alpha$ maakt met de positieve
$x$-as, dan geldt:

\[
\boxed{
A_x=A\cos\alpha
}
\]

en:

\[
\boxed{
A_y=A\sin\alpha.
}
\]


\begin{voorbeeld}[Een kracht ontbinden]

Een kracht heeft grootte:

\[
F=100\,\mathrm N
\]

en maakt een hoek van:

\[
30^\circ
\]

met de horizontale.

De horizontale component is:

\[
F_x
=
F\cos30^\circ.
\]

Dus:

\[
F_x
=
100\cos30^\circ
\approx
86{,}6\,\mathrm N.
\]

De verticale component is:

\[
F_y
=
F\sin30^\circ.
\]

Dus:

\[
F_y
=
100\sin30^\circ
=
50{,}0\,\mathrm N.
\]

Daarmee hebben we de oorspronkelijke kracht vervangen door twee
loodrechte componenten:

\[
\boxed{
\vec F=\vec F_x+\vec F_y.
}
\]

\end{voorbeeld}


\begin{opmerking}

Het ontbinden van vectoren is geen doel op zich.

Het is een techniek die fysische problemen eenvoudiger maakt omdat
we effecten in verschillende richtingen afzonderlijk kunnen
onderzoeken.

In het volgende hoofdstuk zullen we dit voortdurend gebruiken bij
de wetten van Newton.

\end{opmerking}


\begin{oefening}[Vectorcomponenten]{\oefB\ESS}

Een kracht van:

\[
40\,\mathrm N
\]

maakt een hoek van:

\[
60^\circ
\]

met de positieve horizontale as.

\begin{question}
Bereken de horizontale component.
\end{question}

\begin{question}
Bereken de verticale component.
\end{question}

\begin{question}
Maak een verzorgde vectorschets.
\end{question}

\end{oefening}


% ============================================================
% 14 TERUG VAN COMPONENTEN
% ============================================================

\FVDsection{Van componenten naar grootte en richting}

Wanneer de loodrechte componenten gekend zijn, kunnen we de grootte van
de oorspronkelijke vector terugvinden met de stelling van Pythagoras:

\[
\boxed{
A
=
\sqrt{A_x^2+A_y^2}.
}
\]

De richting kan worden bepaald met:

\[
\boxed{
\tan\alpha
=
\frac{A_y}{A_x}.
}
\]


\begin{voorbeeld}

Een vector heeft componenten:

\[
A_x=6{,}0
\]

en:

\[
A_y=8{,}0.
\]

De grootte is:

\[
A
=
\sqrt{6{,}0^2+8{,}0^2}.
\]

Dus:

\[
\boxed{
A=10.
}
\]

Voor de hoek:

\[
\tan\alpha
=
\frac{8}{6}.
\]

Daaruit:

\[
\alpha
\approx
53^\circ.
\]

\end{voorbeeld}


\begin{opmerking}

Let bij het berekenen van een hoek altijd op het kwadrant waarin de
vector ligt.

Alleen:

\[
\tan^{-1}
\left(
\frac{A_y}{A_x}
\right)
\]

berekenen is niet altijd voldoende om de fysisch juiste richting te
vinden.

Maak daarom altijd eerst een schets.

\end{opmerking}


\begin{oefening}[Van componenten naar vector]{\oefB\ESS}

Een vector heeft componenten:

\[
A_x=-3{,}0
\]

en:

\[
A_y=4{,}0.
\]

\begin{question}
Schets de vector in een assenstelsel.
\end{question}

\begin{question}
In welk kwadrant ligt de vector?
\end{question}

\begin{question}
Bereken de grootte.
\end{question}

\begin{question}
Bepaal de richting ten opzichte van de positieve $x$-as.
\end{question}

\end{oefening}


% ============================================================
% 15 GEÏNTEGREERDE VECTOROEFENING
% ============================================================

\FVDsection{Afstand en verplaatsing zijn niet hetzelfde}

Vectoren laten ons onderscheid maken tussen grootheden die op het
eerste gezicht sterk op elkaar lijken.


\begin{voorbeeld}[Een reddingsdrone]

Een reddingsdrone vliegt eerst:

\[
120\,\mathrm m
\]

naar het oosten en daarna:

\[
50\,\mathrm m
\]

naar het noorden.

De afgelegde afstand is:

\[
s
=
120+50
=
170\,\mathrm m.
\]

Afstand is hier een scalaire grootheid.

De verplaatsingsvector heeft componenten:

\[
\Delta x
=
120\,\mathrm m,
\]

\[
\Delta y
=
50\,\mathrm m.
\]

De grootte van de verplaatsing is:

\[
|\Delta\vec r|
=
\sqrt{120^2+50^2}.
\]

Dus:

\[
\boxed{
|\Delta\vec r|
=
130\,\mathrm m.
}
\]

Voor de richting:

\[
\tan\alpha
=
\frac{50}{120}.
\]

Daaruit:

\[
\alpha
\approx
22{,}6^\circ.
\]

De verplaatsing is dus ongeveer:

\[
\boxed{
23^\circ
\text{ ten noorden van het oosten}.
}
\]

\end{voorbeeld}


\begin{oefening}[Robotroute]{\oefU\ESS}

Een robot rijdt eerst:

\[
6{,}0\,\mathrm m
\]

naar het oosten en daarna:

\[
8{,}0\,\mathrm m
\]

naar het noorden.

\begin{question}
Bereken de afgelegde afstand.
\end{question}

\begin{question}
Bereken de grootte van de verplaatsing.
\end{question}

\begin{question}
Bepaal de richting van de verplaatsing.
\end{question}

\begin{question}
Leg uit waarom afstand en grootte van de verplaatsing verschillend
zijn.
\end{question}

\begin{question}
Welke grootheid is scalair en welke vectorieel?
\end{question}

\end{oefening}


% ============================================================
% 16 DIMENSIONELE CONTROLE
% ============================================================

\FVDsection{Eenheden als controlemiddel}

Eenheden zijn niet alleen een toevoeging achter een antwoord.

Ze kunnen helpen om fouten te vinden.

Stel dat iemand schrijft:

\[
v=\frac{s}{t^2}.
\]

De eenheid van het rechterlid is:

\[
\frac{\mathrm m}{\mathrm{s^2}}
=
\mathrm{m/s^2}.
\]

Dat is de eenheid van een versnelling en niet van een snelheid.

De formule kan dus onmogelijk een correcte formule voor snelheid zijn.


\begin{definitie}[Dimensionele consistentie]

Een fysische vergelijking is \textbf{dimensioneel consistent} wanneer
de dimensies van beide leden met elkaar overeenstemmen.

\end{definitie}


\begin{voorbeeld}

Controleer:

\[
s=vt.
\]

Rechts vinden we:

\[
[v][t]
=
\mathrm{m/s}\cdot\mathrm s.
\]

Dus:

\[
[v][t]
=
\mathrm m.
\]

Dat komt overeen met de eenheid van $s$.

De vergelijking is dimensioneel consistent.

\end{voorbeeld}


\begin{opmerking}

Dimensionele consistentie bewijst niet dat een formule fysisch juist
is.

Een dimensioneel inconsistente formule kan echter wel onmiddellijk
worden verworpen.

\end{opmerking}


\begin{oefening}[Controleer de eenheden]{\oefU\ESS}

Onderzoek of de volgende formules dimensioneel mogelijk zijn.

\begin{enumerate}
  \item
  \[
  v=at
  \]

  \item
  \[
  s=vt
  \]

  \item
  \[
  s=at
  \]

  \item
  \[
  E=Fs
  \]

  waarbij $[F]=\mathrm N$ en $[s]=\mathrm m$.

\end{enumerate}

Leg telkens je redenering uit.

\end{oefening}


% ============================================================
% 17 PROBLEEMOPLOSSEN
% ============================================================

\FVDsection{Een fysisch probleem systematisch aanpakken}

Een fysisch probleem oplossen is meer dan een formule zoeken waarin
de juiste letters voorkomen.

Bij veel problemen is de moeilijkste stap:

\[
\boxed{
\text{werkelijke situatie}
\longrightarrow
\text{fysisch model}.
}
\]

Doorheen deze cursus gebruiken we daarom een systematische aanpak.


\FVDsubsection{Stap 1 — Begrijp de situatie}

Lees het probleem volledig.

Vraag jezelf af:

\begin{itemize}
  \item Wat gebeurt er?
  \item Welke voorwerpen of systemen spelen een rol?
  \item Welke informatie is relevant?
  \item Welke informatie is eventueel overbodig?
\end{itemize}


\FVDsubsection{Stap 2 — Maak een schets}

Een eenvoudige schets maakt vaak onmiddellijk duidelijk:

\begin{itemize}
  \item welke richting positief wordt gekozen;
  \item welke afstanden of hoeken belangrijk zijn;
  \item welke vectoren later nodig zullen zijn.
\end{itemize}


\FVDsubsection{Stap 3 — Noteer gegeven en gevraagd}

Gebruik correcte symbolen en eenheden.

Bijvoorbeeld:

\[
\begin{array}{ll}
\text{gegeven:} &
m=2{,}0\,\mathrm{kg},\\
&
a=3{,}0\,\mathrm{m/s^2},
\\[1ex]
\text{gevraagd:} &
F=?
\end{array}
\]


\FVDsubsection{Stap 4 — Kies een fysisch model}

Vraag jezelf af welke fysische principes bruikbaar zijn.

Later kunnen dat bijvoorbeeld zijn:

\[
\sum\vec F=m\vec a,
\]

de gravitatiewet,

\[
F_g
=
G\frac{m_1m_2}{r^2},
\]

of de wet van Coulomb.


\FVDsubsection{Stap 5 — Werk indien mogelijk symbolisch}

Probeer de gevraagde grootheid eerst symbolisch vrij te maken.

Dat maakt zichtbaar hoe het antwoord afhangt van de verschillende
grootheden en verkleint de kans op invoerfouten.


\FVDsubsection{Stap 6 — Reken uit}

Vul pas daarna de getalwaarden in.

Gebruik geschikte eenheden en rond verantwoord af.


\FVDsubsection{Stap 7 — Controleer het resultaat}

Vraag jezelf af:

\begin{itemize}
  \item Heeft het antwoord de juiste eenheid?
  \item Is de grootteorde realistisch?
  \item Is het teken logisch?
  \item Is de richting logisch?
  \item Heb ik verantwoord afgerond?
  \item Beantwoordt mijn resultaat werkelijk de gestelde vraag?
\end{itemize}


\begin{opmerking}

Een fysische berekening eindigt niet wanneer de rekenmachine een getal
geeft.

Een oplossing is pas volledig wanneer het resultaat gecontroleerd en
fysisch geïnterpreteerd is.

\end{opmerking}


% ============================================================
% 18 OEFENINGEN
% ============================================================

\FVDsection{Geïntegreerde oefeningen}


\begin{oefening}[Van nanometer tot kilometer]{\oefB\ESS}

Zet alle waarden om naar SI-eenheden en schrijf ze daarna in
wetenschappelijke notatie.

\begin{enumerate}
  \item $450\,\mathrm{nm}$
  \item $3{,}2\,\mathrm{mm}$
  \item $7{,}5\,\mathrm{km}$
  \item $25\,\mathrm{\mu s}$
  \item $3{,}6\,\mathrm{cm^2}$
\end{enumerate}

\end{oefening}


\begin{oefening}[Dichtheid uit meetgegevens]{\oefB\ESS}

Een laborant meet:

\[
m=0{,}00450\,\mathrm{kg}
\]

en:

\[
V=1{,}20\cdot10^{-6}\,\mathrm{m^3}.
\]

De dichtheid wordt gegeven door:

\[
\rho=\frac{m}{V}.
\]

\begin{question}
Bereken de dichtheid.
\end{question}

\begin{question}
Rapporteer het resultaat met een verantwoord aantal beduidende cijfers.
\end{question}

\begin{question}
Controleer de eenheid.
\end{question}

\end{oefening}


\begin{oefening}[Welke vector hoort erbij?]{\oefB}

Een drone beweegt met:

\[
12\,\mathrm{m/s}
\]

naar het westen.

\begin{question}
Welke informatie bevat alleen het getal $12\,\mathrm{m/s}$?
\end{question}

\begin{question}
Waarom is die informatie onvoldoende om de snelheid als vector volledig
te beschrijven?
\end{question}

\begin{question}
Maak een vectorschets.
\end{question}

\end{oefening}


\begin{oefening}[Twee krachten combineren]{\oefU\ESS}

Op een voorwerp werken twee loodrechte krachten:

\[
F_x=30\,\mathrm N
\]

naar rechts en:

\[
F_y=40\,\mathrm N
\]

naar boven.

\begin{question}
Teken beide vectoren.
\end{question}

\begin{question}
Bepaal grafisch de richting van de resultante.
\end{question}

\begin{question}
Bereken de grootte van de resultante.
\end{question}

\begin{question}
Bereken de hoek van de resultante met de positieve horizontale richting.
\end{question}

\end{oefening}


\begin{oefening}[Een slecht experiment]{\oefU\ESS}

Een leerling wil onderzoeken hoe de lengte van een metalen draad
afhangt van zijn temperatuur.

Hij meet de lengte bij verschillende temperaturen, maar gebruikt bij
iedere temperatuur een andere draad van hetzelfde materiaal.

\begin{question}
Wat is de onafhankelijke variabele?
\end{question}

\begin{question}
Wat is de afhankelijke variabele?
\end{question}

\begin{question}
Waarom is het gebruiken van telkens een andere draad problematisch?
\end{question}

\begin{question}
Hoe zou je het experiment verbeteren?
\end{question}

\end{oefening}


\begin{oefening}[Ontwerp een experiment]{\oefG}

Een onderzoeker wil nagaan of de uitrekking van een elastiek recht
evenredig is met de aangelegde kracht.

Ontwerp een experiment.

Beschrijf:

\begin{enumerate}
  \item de onderzoeksvraag;
  \item een toetsbare hypothese;
  \item de onafhankelijke variabele;
  \item de afhankelijke variabele;
  \item minstens drie controlevariabelen;
  \item welke grootheden worden gemeten;
  \item welke eenheden worden gebruikt;
  \item hoe de meetgegevens in een grafiek worden voorgesteld;
  \item welk grafisch patroon je verwacht bij recht evenredigheid;
  \item hoe je beoordeelt of de meetgegevens voldoende bij het model passen;
  \item minstens twee mogelijke bronnen van meetonzekerheid.
\end{enumerate}

\end{oefening}


\begin{oefening}[Navigeren met wind]{\oefG}

Een drone moet rechtstreeks van punt $A$ naar punt $B$ vliegen.

Punt $B$ ligt ten opzichte van $A$:

\[
240\,\mathrm m
\]

naar het oosten en:

\[
180\,\mathrm m
\]

naar het noorden.

De drone kan ten opzichte van de lucht vliegen met een snelheid van:

\[
15{,}0\,\mathrm{m/s}.
\]

Er staat een constante wind met snelheid:

\[
4{,}0\,\mathrm{m/s}
\]

naar het oosten.

Onderzoek in welke richting de drone ten opzichte van de lucht moet
vliegen om rechtstreeks van $A$ naar $B$ te bewegen.

Maak eerst een duidelijke vectorschets en beschrijf je
oplossingsstrategie voordat je begint te rekenen.

\end{oefening}


% ============================================================
% 19 VEELVOORKOMENDE FOUTEN
% ============================================================

\FVDsection{Veelvoorkomende denkfouten}

\begin{itemize}

  \item \textbf{Een getal zonder eenheid noteren.}

  In een fysische context is:

  \[
  12
  \]

  meestal onvoldoende.

  \item \textbf{Alle cijfers van de rekenmachine overnemen.}

  De rekenmachine kent de nauwkeurigheid van je metingen niet.

  \item \textbf{Het aantal decimalen verwarren met het aantal
  beduidende cijfers.}

  Bijvoorbeeld:

  \[
  0{,}00450
  \]

  heeft drie beduidende cijfers.

  \item \textbf{Een oppervlakte-eenheid lineair omzetten.}

  Uit:

  \[
  1\,\mathrm{cm}=10^{-2}\,\mathrm m
  \]

  volgt:

  \[
  1\,\mathrm{cm^2}=10^{-4}\,\mathrm{m^2}.
  \]

  \item \textbf{Vectoren behandelen alsof ze gewone positieve
  getallen zijn.}

  Richting en zin zijn essentieel.

  \item \textbf{Een berekening uitvoeren zonder eerst een schets of
  model te maken.}

  Bij complexere problemen leidt dit vaak tot teken- of
  interpretatiefouten.

  \item \textbf{Een afwijkende meting zomaar verwijderen.}

  Eerst moet onderzocht worden waarom die meting afwijkt.

\end{itemize}


% ============================================================
% 20 SAMENVATTING
% ============================================================

\FVDsection{Samenvatting}

Een fysische grootheid wordt beschreven met een getalwaarde en eenheid:

\[
\boxed{
\text{grootheid}
=
\text{getalwaarde}
\times
\text{eenheid}.
}
\]

Voor wetenschappelijke communicatie gebruiken we bij voorkeur
SI-eenheden.

Zeer grote en zeer kleine waarden kunnen worden geschreven als:

\[
\boxed{
a\cdot10^n
}
\]

met:

\[
1\leq |a|<10.
\]

Een meetresultaat heeft een beperkte nauwkeurigheid.

Het aantal gerapporteerde cijfers moet daarmee overeenkomen.

Bij vermenigvuldigen en delen bepaalt in deze cursus doorgaans de
meetwaarde met het kleinste aantal beduidende cijfers de rapportering
van het eindresultaat.

Bij optellen en aftrekken kijken we naar de plaats van het laatste
betrouwbare cijfer.

Een scalaire grootheid wordt door haar grootte beschreven.

Een vectoriële grootheid heeft:

\[
\boxed{
\text{grootte, richting en zin}.
}
\]

Een vector kan worden ontbonden:

\[
\boxed{
A_x=A\cos\alpha
}
\]

en:

\[
\boxed{
A_y=A\sin\alpha.
}
\]

Uit componenten volgt:

\[
\boxed{
A=\sqrt{A_x^2+A_y^2}.
}
\]

De richting kan worden bepaald met:

\[
\boxed{
\tan\alpha=\frac{A_y}{A_x}.
}
\]

Een fysische vergelijking moet dimensioneel consistent zijn.

Bij experimenteel onderzoek onderscheiden we:

\[
\boxed{
\text{onafhankelijke variabele},
\quad
\text{afhankelijke variabele},
\quad
\text{controlevariabelen}.
}
\]

Een goede probleemaanpak verloopt bijvoorbeeld als:

\[
\boxed{
\text{begrijpen}
\rightarrow
\text{schetsen}
\rightarrow
\text{gegevens}
\rightarrow
\text{model}
\rightarrow
\text{rekenen}
\rightarrow
\text{controleren}
\rightarrow
\text{interpreteren}.
}
\]


% ============================================================
% 21 MINIMALE BEHEERSING
% ============================================================

\FVDsection{Wat moet je zeker beheersen?}

Voor je verdergaat naar de wetten van Newton moet je de onderstaande
vaardigheden voldoende beheersen.

Je moet:

\begin{itemize}

  \item fysische grootheden en eenheden correct kunnen noteren;

  \item courante SI-voorvoegsels kunnen gebruiken;

  \item eenheden kunnen omzetten;

  \item ook oppervlakte- en volume-eenheden correct kunnen omzetten;

  \item wetenschappelijke notatie kunnen lezen en schrijven;

  \item het aantal beduidende cijfers van een meetwaarde kunnen bepalen;

  \item een berekend resultaat verantwoord kunnen afronden;

  \item een meetresultaat kritisch kunnen beoordelen;

  \item scalaire en vectoriële grootheden kunnen onderscheiden;

  \item een vector correct kunnen voorstellen;

  \item vectoren kunnen optellen;

  \item een vector kunnen ontbinden in horizontale en verticale
        componenten;

  \item uit componenten grootte en richting kunnen terugvinden;

  \item eenheden kunnen gebruiken als controle op een fysische formule;

  \item onafhankelijke, afhankelijke en controlevariabelen kunnen
        onderscheiden;

  \item een fysisch probleem systematisch kunnen aanpakken.

\end{itemize}

De oefeningen met de markering \verb|\ESS| zijn noodzakelijk om deze
gereedschapskist voldoende te beheersen.

Niet alle essentiële oefeningen zijn \verb|\oefB|-oefeningen.
Sommige \verb|\oefU|-oefeningen zijn eveneens noodzakelijk omdat je
niet alleen technieken moet uitvoeren, maar meetgegevens, modellen en
redeneringen ook kritisch moet kunnen analyseren.


% ============================================================
% 22 VOORUITBLIK
% ============================================================

\FVDsection{Vooruitblik: van vectoren naar krachten}

We beschikken nu over een aantal belangrijke gereedschappen.

We kunnen:

\[
\text{meten},
\]

\[
\text{eenheden gebruiken},
\]

\[
\text{vectoren voorstellen},
\]

\[
\text{vectoren ontbinden},
\]

en:

\[
\text{fysische problemen systematisch analyseren}.
\]

In het volgende hoofdstuk gebruiken we deze gereedschappen voor een
fundamentele vraag:

\begin{center}
\textit{Waarom verandert de beweging van een voorwerp?}
\end{center}

Daarvoor onderzoeken we krachten en de drie wetten van Newton.

De stap wordt:

\[
\boxed{
\text{vectoren}
\longrightarrow
\text{krachten}
\longrightarrow
\sum\vec F
\longrightarrow
\text{bewegingsverandering}.
}
\]

\end{document}